% Part: computability
% Chapter: recursive-functions
% Section: normal-form

\documentclass[../../../include/open-logic-section]{subfiles}

\begin{document}

\olfileid{cmp}{rec}{nft}
\olsection{قضیهٔ صورت نرمال}

\begin{thm}[قضیهٔ صورت نرمال کلینی]
\ollabel{thm:kleene-nf}
رابطه‌ای بازگشتی اولیه مانند $T(e,x,s)$ و تابعی بازگشتی اولیه مانند
$U(s)$ وجود دارند که ویژگی زیر را دارند: اگر $f$ هر تابع بازگشتی
جزئی باشد، آنگاه برای یک~$e$،
\[
f(x) \simeq U(\umin{s}{T(e, x, s)})
\]
به ازای هر $x$ برقرار است.
\end{thm}

\begin{explain}
برهان قضیهٔ صورت نرمال مفصل است، اما اندیشهٔ اصلی ساده است. هر تابع
بازگشتی جزئی یک \emph{نمایه}~$e$ دارد؛ به‌طور شهودی، عددی که برنامه
یا تعریف آن را کد می‌کند. اگر $f(x)\fdefined$، می‌توان محاسبه را
به‌طور نظام‌مند ثبت کرد و با عددی چون~$s$ کد کرد، و این واقعیت که
$s$ محاسبهٔ~$f$ را برای ورودی~$x$ کد می‌کند تنها با استفاده از $x$
و تعریف~$e$ به‌صورت بازگشتی اولیه قابل بررسی است. بنابراین رابطهٔ
$T$، یعنی «تابع دارای نمایهٔ~$e$ برای ورودی~$x$ محاسبه‌ای دارد و
$s$ این محاسبه را کد می‌کند»، بازگشتی اولیه است. با داشتن ثبت کامل
محاسبه، یعنی~$s$، «حاصل» آن مقدار~$f(x)$ است و آن را نیز می‌توان
به‌صورت بازگشتی اولیه از~$s$ به دست آورد.

قضیهٔ صورت نرمال نشان می‌دهد که برای تعریف هر تابع بازگشتی جزئی تنها
یک جست‌وجوی نامحدود لازم است. در اصل می‌توان در همهٔ اعداد جست‌وجو
کرد تا عددی بیابیم که محاسبهٔ تابعِ دارای نمایهٔ~$e$ را برای
ورودی~$x$ کد می‌کند. می‌توانیم عددهای~$e$ را «نام» توابع بازگشتی
جزئی قرار دهیم و $\cfind{e}$ را برای تابع~$f$ای بنویسیم که معادلهٔ
قضیه آن را تعریف می‌کند. توجه کنید که هر تابع بازگشتی جزئی ممکن است
بیش از یک نمایه داشته باشد---در واقع هر تابع بازگشتی جزئی بی‌نهایت
نمایه دارد.
\end{explain}
\end{document}
